\section{Accounting, Controls and Reported Performance}\label{sec:accounting}
% Generated by code/verify_part2_accounting.py; do not edit. Invoice example source: data/examples/invoice_chain.json.
\providecommand{\acctExScorepathG}{0.17}
\providecommand{\acctExCashOne}{0.32}
\providecommand{\acctExRevenue}{\$1.2 billion}
\providecommand{\acctExTpar}{0.5}
\providecommand{\acctExBuilds}{\$622{,}000}
\providecommand{\acctExIntegration}{\$2.00 million}
\providecommand{\acctExUnc}{40\%}
\providecommand{\acctExDC}{\$500{,}000}
\providecommand{\acctExFcap}{\$1.87 million}
\providecommand{\acctExFcapShare}{0.67}
\providecommand{\acctExTB}{4}
\providecommand{\acctExGone}{0.21}
\providecommand{\acctExGfull}{$-0.05$}
\providecommand{\acctExFullEnd}{4}
\providecommand{\acctExPartialYear}{5}
\providecommand{\acctExGpartial}{$-0.03$}
\providecommand{\acctExCashLater}{$-0.01$}
\providecommand{\acctExOther}{\$186{,}000}
\providecommand{\acctExSevShift}{0.005}
\providecommand{\acctExGoneSev}{0.21}
\providecommand{\acctTaxAfter}{0.25}
\providecommand{\acctTaxForeign}{0.011}
\providecommand{\acctTaxPrior}{0.033}
\providecommand{\acctTaxSD}{\$2.12 million}
\providecommand{\acctCutOpen}{58\%}
\providecommand{\acctCutWait}{0.087}
\providecommand{\acctCutCost}{\$43{,}000}
\providecommand{\acctCutDelay}{\$500{,}000}
\providecommand{\acctCutPafter}{2.4\%}
\providecommand{\acctCutPlate}{9\%}
\providecommand{\acctCutZstar}{50\%}
\providecommand{\acctAminLine}{\$205{,}000}
\providecommand{\acctAminFiveCost}{\$755{,}000}
\providecommand{\acctAminOneCost}{\$805{,}000}
\providecommand{\acctAminFiveMiss}{3.1\%}
\providecommand{\acctAminOneMiss}{0.7\%}
\providecommand{\acctAminFiveResid}{\$310{,}000}
\providecommand{\acctAminOneResid}{\$67{,}000}
\providecommand{\acctAminFloorAI}{11.3\%}
\providecommand{\acctAminFloorHu}{4.7\%}
\providecommand{\acctAminHumThreeMiss}{5.0\%}
\providecommand{\acctAminHumThreeCost}{\$1.50 million}
\providecommand{\acctAminLossCost}{\$872{,}000}
\providecommand{\acctAminWeakOneCost}{\$1.25 million}
\providecommand{\acctAminItgcValue}{\$450{,}000}
\providecommand{\acctAminPbarFive}{0.00213}
\providecommand{\acctAminPbarOne}{0.00042}
\providecommand{\acctItgcInv}{\$205{,}000}
\providecommand{\acctItgcRef}{\$75{,}000}
\providecommand{\acctItgcUp}{\$20{,}000}
\providecommand{\acctItgcIfc}{\$40{,}000}
\providecommand{\acctItgcNint}{4}
\providecommand{\acctItgcTup}{1.22}
\providecommand{\acctItgcPVup}{\$57{,}000}
\providecommand{\acctItgcPVflow}{\$61{,}000}
\providecommand{\acctItgcNPV}{\$0.52 million}
\providecommand{\acctItgcNPVafter}{$-$\$4.30 million}
\providecommand{\acctItgcInvEng}{\$205{,}000}
\providecommand{\acctItgcRefEng}{\$75{,}000}
\providecommand{\acctItgcApp}{\$88{,}000}
\providecommand{\acctItgcAppEng}{\$88{,}000}
\providecommand{\acctItgcVal}{\$24{,}000}
\providecommand{\acctItgcValEng}{\$24{,}000}
\providecommand{\acctItgcNPVEng}{\$0.52 million}
\providecommand{\acctItgcPos}{0.56}
\providecommand{\acctItgcShare}{0.020\%}
\providecommand{\acctItgcBandNFive}{8}
\providecommand{\acctItgcBandNOne}{7}
\providecommand{\acctItgcBandNfbFive}{4}
\providecommand{\acctItgcBandNLo}{13}
\providecommand{\acctItgcFloorFive}{0.37\%}
\providecommand{\acctItgcFloorOne}{0.38\%}
\providecommand{\acctItgcFloorOld}{0.22\%}
\providecommand{\acctAminDirect}{0.30\%}
\providecommand{\acctAminDirectLo}{0.18\%}
\providecommand{\acctAminDirectOld}{0.14\%}
\providecommand{\acctAminKpmgSmall}{\$1.5 million}
\providecommand{\acctAminAuditFee}{\$1.49 million}
\providecommand{\acctAminPccFive}{0.00084}
\providecommand{\acctAminPccOne}{0.00014}
\providecommand{\acctAminPhmFive}{8.4\%}
\providecommand{\acctAminSepOne}{\$1.01 million}
\providecommand{\acctItgcTests}{10}
% example numbers, generated by code/verify_part2_accounting.py from data/examples/invoice_chain.json

This section replaces the stylized GAAP basis of Section~\ref{sec:tasks} with U.S.\ standards and tax law, and adds the constraint that internal control over financial reporting (ICFR) places on automating financial processes. Capitalization never changes the cash path but makes the GAAP dip shallower and longer; a growing program keeps GAAP above cash permanently; and under ASU 2025-06 a novel AI build is expensed until it is shown to work. Where a build is performed changes its after-tax cost by about a tenth. Control redesign is a cost and a calendar constraint. Because ICFR is pass or fail, the accountability floor $a_{\min}$ in Corollary~\ref{cor:ceiling} is the cost of the cheapest control design that leaves no material weakness, and AI-run controls bring a common-mode risk that more layers on one model cannot remove, so some cheap designs are not available at any price and $a_{\min}$ does not vanish and can rise. Automating a step never takes it out of ICFR scope: the agent becomes an in-scope application with its own IT general controls, retested at every forced model upgrade, and provider outages in the close window put the reporting calendar at risk. Law and standards are as of 29 October 2026; where guidance is silent, our treatments are judgments. Proofs are in Appendix~\ref{app:proofs-accounting}, checks in \texttt{code/verify\_part2\_accounting.py}.

\subsection{What the standards say}\label{sec:acct-standards}

\paragraph{Internal-use software.} FASB ASC Subtopic 350-40 \citep{fasb2026asc35040} capitalizes costs only in the application development stage of a sequence of project stages; ASU 2025-06 (September 2025) removes the stages \citep{fasb2025asu202506}. Capitalization begins when management, with the relevant authority, authorizes and commits to funding the project and it is probable that the project will be completed and the software used to perform the function intended (ASC 350-40-25-12). The threshold is not met while significant development uncertainty exists: technological innovations or novel, unique or unproven functions or features not yet resolved through coding and testing, or significant performance requirements not yet identified or still being substantially revised (25-12A). Eligible costs are unchanged: external direct costs of materials and services consumed, payroll costs of employees directly associated with the project, and interest (30-1). Data conversion is expensed, except software to access or convert old data (25-5, 30-1(d)), and so is training (25-4). Capitalization ceases no later than when the software is substantially complete and ready for its intended use (25-14); amortization, generally straight line over a life reflecting obsolescence, technology and competition, begins when it is ready for its intended use (35-4 to 35-6). The Update is effective for annual periods beginning after 15 December 2027; early adoption is permitted (65-4).

\paragraph{Hosted arrangements.} Implementation costs of a vendor's hosted software under a service contract are capitalized under the same guidance (ASU 2018-15; 25-18), expensed over the hosting term including renewals reasonably certain to be exercised, and presented with the hosting fees in the income statement, balance sheet and cash-flow statement (35-13, 35-14, 45-1 to 45-3) \citep{fasb2018asu201815}: they are operating outflows, whereas capitalized internal-use software is investing.

\paragraph{What the guidance does not address.} Neither ASC 350-40 nor ASU 2025-06 mentions model training, fine-tuning, prompt and agent development, evaluation datasets or data labeling. A practitioner illustration expenses a generative-AI feature's costs until coding and testing resolve the uncertainty about its novel function, then capitalizes the model training that follows, and treats purchased training data with an alternative use as an intangible under Subtopic 350-30 \citep{deloitte2025headsup}; whether data cleansing and labeling are data conversion is unaddressed \citep{deloitte2026aicosts}. Table~\ref{tab:acct-map} flags these as judgments.

\paragraph{Exit costs and abandoned assets.} Under ASC Topic 420 \citep{fasb2026asc420}, a one-time termination benefit is recognized at the communication date if employees need not serve beyond the minimum retention period (at most the legal notification period, or 60 days), and otherwise ratably over the future service period (420-10-25-7 to 25-9). Benefits under an ongoing arrangement or past practice fall under Topic 712 and accrue when payment is probable and reasonably estimable \citep{fasb2026asc712}, which can precede communication. Early contract termination costs (leases excluded) are recognized when the entity terminates under the contract's terms (420-10-25-12). Legacy assets to be abandoned are not written off at cutover: from the commitment to abandon, they are depreciated over the shortened useful life, down to salvage value at the cease-use date (ASC 360-10-35-47, 35-48) \citep{fasb2026asc360}. Previously recognized share-based compensation cost is reversed for awards that displaced employees forfeit \citep{fasb2016asu201609}.

\paragraph{Tax.} Domestic research or experimental expenditures paid or incurred in taxable years beginning after 31 December 2024 are deductible when paid or incurred, or by election over at least 60 months (26 U.S.C.~\S174A(a), (c)); foreign ones are amortized ratably over 15 years from the midpoint of the taxable year, with no deduction on abandonment (\S174(a), (d)); and software development costs are research or experimental expenditures (\S174A(d)(3), \S174(c)(3)) \citep{usc2026sec174a,usc2026sec174}. Rev.\ Proc.\ 2025-28 sets the transition elections \citep{irs2025revproc202528}. Interim guidance under the earlier rules excludes configuring and deploying purchased software, training, routine maintenance and data conversion from software development \citep{irs2023notice202363}; whether it carries over to \S174A awaits guidance. Qualified property acquired after 19 January 2025, which includes computer equipment such as servers (five-year property), gets a permanent 100\% first-year depreciation allowance under \S168(k) as amended by \S70301 of Pub.\ L.\ 119-21; a taxpayer can elect out for a class of property (\S168(k)(7)), and property used predominantly outside the United States is depreciated under the alternative depreciation system and does not qualify \citep{irs2026notice202611,usc2026sec168,cfr2026sec1168k2}. Book and tax treatments are independent; differences are temporary differences under ASC Topic 740 \citep{fasb2026asc740}.

\begin{table}[t]
\centering
\caption{Transformation costs: GAAP treatment, timing of P\&L and cash, and U.S.\ tax treatment}\label{tab:acct-map}
\scriptsize
\setlength{\tabcolsep}{3pt}
\begin{tabular}{@{}p{2.6cm}p{5.0cm}p{2.5cm}p{3.6cm}@{}}
\toprule
Cost (Section~\ref{sec:tasks}) & GAAP & Cash & U.S.\ tax\\
\midrule
In-house build ($F$): coding, configuration, testing & ASC 350-40: capitalize once 25-12 is met (stage model before adoption); amortize from ready for intended use & Investing & Deduct if domestic (\S174A); 15 years if foreign (\S174)\\
Model training, fine-tuning, evaluation suites & Silent. Judgment: development cost, expensed while significant development uncertainty exists & Investing if capitalized & Likely software development; judgment\\
Prompt and agent work on a vendor platform & Hosted implementation (25-18): capitalize as above; expense over hosting term in the fee line & Operating & Not development under Notice 2023-63; judgment for \S174A\\
Data conversion, cleansing, labeling; staff training; purchased data & Expense (25-4, 25-5), except software to access or convert old data; labeling: judgment; purchased data with alternative use: intangible (350-30) & Operating; investing if capitalized & Not development (Notice 2023-63)\\
Equipment ($I$) & PP\&E (ASC 360), depreciated & Investing & 100\% first-year allowance\\
Parallel run, dip, control redesign ($C^{\mathrm{par}},C^J,C^{\mathrm{ctl}}$) & Expense as incurred & Operating & Deduct\\
Severance, break fees ($C^{\mathrm{sev}},C^{\mathrm{brk}}$) & ASC 712 plan: when probable and estimable; ASC 420: at communication, or ratably over required service; break fees at termination under contract terms & Operating, when paid & Deduct; timing difference\\
Legacy assets ($C^{\mathrm{wo}}$) & Depreciate over shortened life, commitment to cease-use (360-10-35-47); impair if indicated & None & Foreign R\&E: no deduction on abandonment\\
Run costs, revalidation; SBC of displaced staff & Expense (an upgrade adding functionality is capitalized); ASC 718: reverse cost of forfeited awards & Operating; none for SBC & Deduct, or R\&E\\
\bottomrule
\end{tabular}

\smallskip
\begin{minipage}{0.98\linewidth}\scriptsize
Notes: references are to FASB ASC 350-40 as amended by ASU 2025-06 unless stated, and to 26 U.S.C.\ as amended by Pub.\ L.\ 119-21. ``Judgment'' marks entries where the guidance is silent; the engine exposes each as a policy input.
\end{minipage}
\end{table}

\subsection{Capitalization and the score path}\label{sec:acct-path}

\begin{assumption}[Books for a transition]\label{ass:acct-books}
Take Assumption~\ref{ass:tasks-transition} on its success path with realized share $b$, fiscal years $[y-1,y)$ for $y\in\mathbb Z$, and equipment $I$ depreciated as in Proposition~\ref{prop:tasks-scorepath}(b).
\begin{enumerate}
\item[(i)] $F=F^{\mathrm{cap}}+F^{\mathrm{exp}}$; $F^{\mathrm{exp}}$ is charged at date 0, and $F^{\mathrm{cap}}$ is amortized straight line over the book life $T_B\ge 1$ from the go-live date $t_{\mathrm{go}}\in[0,T_{\mathrm{par}}]$.
\item[(ii)] From the commitment at date 0, the legacy assets' book value net of salvage in excess of their normal depreciation, $C^{\mathrm{wo}}$, is charged evenly on $[0,T_{\mathrm{par}})$.
\item[(iii)] $C^{\mathrm{str}}=C^{\mathrm{sev}}+C^{\mathrm{brk}}$ is paid at cutover and recognized in GAAP at dates $t_{\mathrm{sev}},t_{\mathrm{brk}}\le T_{\mathrm{par}}$ (possibly negative), or for severance ratably over $[t_{\mathrm{com}},T_{\mathrm{par}})$.
\item[(iv)] The cash score deducts $F$ and $I$ when paid, whatever their classification, as Part I's free cash flow does.
\end{enumerate}
\end{assumption}

\begin{proposition}[Capitalization and the score path]\label{prop:acct-path}
Under Assumption~\ref{ass:acct-books}, let $\Delta R^c_y$ and $\Delta R^G_y$ be the changes in the cash and GAAP scores in fiscal year $y$.
\begin{enumerate}
\item[(a)] (Invariance.) $\Delta R^c_y$ is the cash path of Proposition~\ref{prop:tasks-scorepath}(a) aggregated by year; it does not depend on $F^{\mathrm{cap}}$, $T_B$, $t_{\mathrm{go}}$, recognition dates or classification.
\item[(b)] (Accrual identity.) For every $Y\ge 1$,
\begin{equation}\label{eq:acct-accrual}
S\sum_{y\le Y}\bigl(\Delta R^G_y-\Delta R^c_y\bigr)=I+F^{\mathrm{cap}}(1-\mathsf A_Y)-\mathsf L_Y-C^{\mathrm{wo}}\min\{Y/T_{\mathrm{par}},1\},
\end{equation}
with $\mathsf A_Y$ the share of $F^{\mathrm{cap}}$ amortized and $\mathsf L_Y$ the restructuring charges recognized but unpaid by the end of year $Y$. Eventually the gap is $(I-C^{\mathrm{wo}})/S$: book choices only move charges between years.
\item[(c)] (First year.) If $T_{\mathrm{par}}\le 1\le T_B$, the ramp is linear and the restructuring is recognized in year 1,
\begin{equation}\label{eq:acct-first}
\Delta R^G_1=\Delta R^c_1+\bigl[I+F^{\mathrm{cap}}\bigl(1-(1-t_{\mathrm{go}})/T_B\bigr)-C^{\mathrm{wo}}\bigr]/S ,
\end{equation}
which is Proposition~\ref{prop:tasks-scorepath}(c) when $F^{\mathrm{cap}}=F$, $t_{\mathrm{go}}=0$ and $T_B=T_A$. The gap rises with $F^{\mathrm{cap}}$, by $1-(1-t_{\mathrm{go}})/T_B$ per dollar, and with $t_{\mathrm{go}}$ and $T_B$. If $C^{\mathrm{sev}}$ is instead recognized in year 0, the year-0 score falls by $C^{\mathrm{sev}}/S$ and the year-1 gap rises by as much.
\item[(d)] (Later years.) Under the conditions of (c), for $y\ge 2$, $\Delta R^c_y=bW/S$, and $\Delta R^G_y$ is below it by $F^{\mathrm{cap}}/(T_BS)$ times the part of year $y$ in $[t_{\mathrm{go}},t_{\mathrm{go}}+T_B)$. GAAP reaches its new level only at $t_{\mathrm{go}}+T_B$.
\item[(e)] (ASU 2025-06.) If $F^{\mathrm{dc}}$ (data conversion, training) is expensed and the share $u$ of the rest is incurred under significant development uncertainty, then $F^{\mathrm{cap}}\le(1-u)(F-F^{\mathrm{dc}})$, so the first-year gap is at most $[I+(1-u)(F-F^{\mathrm{dc}})(1-(1-t_{\mathrm{go}})/T_B)-C^{\mathrm{wo}}]/S$, tending to $(I-C^{\mathrm{wo}})/S$ as $u\to 1$.
\end{enumerate}
\end{proposition}

\noindent\emph{Intuition and example.} The GAAP score is the cash score plus the change in net operating assets. In the invoice transition of Section~\ref{sec:tasks} (\acctExRevenue{} of revenue, $T_{\mathrm{par}}=\acctExTpar$ and $b=1$, as in the example of Proposition~\ref{prop:tasks-scorepath}), let \acctExUnc{} of the \acctExBuilds{} of agent builds precede proof that the novel exception-handling function works, let \acctExDC{} of the \acctExIntegration{} integration budget be data conversion and cleansing, and let the \acctExOther{} of redesign, legal review, evaluation data, change management, training and control redesign be expensed (illustrative judgments, except training, 25-4): \acctExFcap{}, a share \acctExFcapShare{} of the build, is capitalized over $T_B=T_A=\acctExTB$ years from cutover. The first-year GAAP score falls by \acctExGone{} points, against \acctExScorepathG{} in Proposition~\ref{prop:tasks-scorepath} and \acctExCashOne{} in cash; it then changes by \acctExGfull{} points a year in years 2 to \acctExFullEnd{} and by \acctExGpartial{} in year \acctExPartialYear, against \acctExCashLater{} in cash, which is negative because the savings net of run costs are. Under an ongoing ASC 712 plan whose terminations become probable before the year begins, \acctExSevShift{} points move into the prior year.

\begin{proposition}[The capitalization premium of a growing program]\label{prop:acct-program}
A firm capitalizes development spending at the rate $C^{\mathrm{cap}}(t)=C^{\mathrm{cap}}_0e^{g_Ft}$, amortizes each vintage straight line over $T_B$ years starting $L\ge 0$ years after spending, and the cash score deducts spending when paid. Let $\mathsf A(x):=(1-e^{-x})/x$, $\mathsf A(0)=1$.
\begin{enumerate}
\item[(a)] Amortization at $t$ is $C^{\mathrm{cap}}(t)e^{-g_FL}\mathsf A(g_FT_B)$, so the GAAP margin exceeds the cash margin by
\begin{equation}\label{eq:acct-wedge}
w(t)=\frac{C^{\mathrm{cap}}(t)}{S(t)}\Bigl[1-e^{-g_FL}\,\mathsf A(g_FT_B)\Bigr].
\end{equation}
\item[(b)] $w$ has the sign of $g_F$. For $g_F>0$ it rises with $g_F$, $T_B$ and $L$, is below $C^{\mathrm{cap}}/S$, and is close to $(C^{\mathrm{cap}}/S)g_F(L+T_B/2)$ when $g_F(L+T_B)$ is small.
\item[(c)] If capitalized spending is a constant share $c^{\mathrm{cap}}$ of revenue growing at $g$, the GAAP Rule-of-40 score exceeds the cash score by $c^{\mathrm{cap}}[1-e^{-gL}\mathsf A(gT_B)]$ every year.
\end{enumerate}
\end{proposition}

\noindent\emph{Intuition and example.} Amortization charges today's income with past, smaller vintages, so the premium is permanent and largest for fast growers. With capitalized development at 5\% of revenue, a three-year life and a half-year lag (illustrative), it is 0.47 points at 5\% growth, 1.26 at 15\% and 2.16 at 30\%: one of the measurement gaps of Section~\ref{sec:discussion}, and a reason for Part I's free cash flow after capitalized software.

\begin{remark}[EBITDA and build or buy]\label{rem:acct-ebitda}
EBITDA adds back amortization of capitalized internal-use software but not of hosted implementation costs, which sit in the fee line (ASC 350-40-45-1), so it favors an in-house build over an equivalent configured platform by the whole capitalized build, as does any cash measure that deducts only purchases of equipment. Part I's free cash flow treats them alike; Proposition~\ref{prop:tasks-buildbuy} should be applied on that basis.
\end{remark}

\begin{proposition}[Tax timing and the after-tax dip]\label{prop:acct-tax}
Let profits be taxed at $\tau$ with full loss offset (Remark~\ref{rem:taxes}); let the build be research or experimental expenditure, $F^{\mathrm{for}}$ of it foreign; let $I$ take the 100\% allowance and other transition costs be deductible when paid. Let $\psi_T$ be the present value per dollar of deductions taken ratably over $T$ years from the midpoint of the first tax year, discounted annually at $r$.
\begin{enumerate}
\item[(a)] $\psi_T=[1+(1+r)^{-T}]/(2T)+[1-(1+r)^{-(T-1)}]/(rT)$, decreasing in $T$ and $r$; at $r=0.10$, $\psi_5=0.796$ and $\psi_{15}=0.532$.
\item[(b)] Under the conditions of Proposition~\ref{prop:acct-path}(c), the after-tax first-year cash score changes by $\Delta R^{c,\tau}_1=(1-\tau)\Delta R^c_1-\tau F^{\mathrm{for}}(1-1/30)/S$.
\item[(c)] A foreign build costs $\tau(1-\psi_{15})$ more per dollar after tax than a domestic one: 9.8 cents at $\tau=0.21$, $r=0.10$.
\item[(d)] The after-tax cash score does not depend on the book treatment; in Part I's notation, $\psi=1$ for domestic software development and fully allowed property, including domestic servers unless the firm elects out, and $\psi=\psi_{15}$ for foreign research.
\end{enumerate}
\end{proposition}

\noindent\emph{Intuition and example.} At $\tau=0.21$ the invoice transition's first-year cash dip of \acctExCashOne{} points becomes \acctTaxAfter{} after tax; building the \acctExBuilds{} of agents abroad deepens it by \acctTaxForeign, and the five-year amortization of 2022--2024, applied to the \acctTaxSD{} of software development (builds, and integration net of data conversion), would have deepened it by \acctTaxPrior. The tax penalty on AI capital in Section~\ref{sec:ext-setup} was 5.7\% under the five-year rule of 2022--2024; under current law it is zero for domestic development and about 12\% for foreign development.

\subsection{Internal control over financial reporting when processes are automated}\label{sec:acct-icfr}

Under Section 404(a) of the Sarbanes--Oxley Act and the SEC's rules, each annual report contains management's assessment of ICFR as of fiscal year-end, and under 404(b) the auditor attests to it, except for emerging growth companies and issuers that are neither accelerated nor large accelerated filers \citep{usc2026sec7262,cfr2026sec229308}. Management may not conclude that ICFR is effective if a material weakness exists, that is, a deficiency, or combination of deficiencies, such that there is a reasonable possibility that a material misstatement of the annual or interim statements will not be prevented or detected on a timely basis \citep{cfr2026sec24012b2}. Changes during the last fiscal quarter that have materially affected, or are reasonably likely to materially affect, ICFR must be disclosed (Item 308(c)); SEC staff cite a new information system as an example and hold management responsible for controls over outsourced operations, and the principal officers certify each periodic report \citep{sec2007faq,usc2026sec7241}. Relief exists for a newly public company's first annual report and, for a year, an acquired business \citep{cfr2026sec229308,sec2007faq}, but not, as far as we found, for a system implementation. Under AS 2201 \citep{pcaob2026as2201}, evidence rises with a control's risk (.46), whose factors include changes in transactions, reliance on IT general controls (ITGCs), automation, and complexity (.47); superseded controls need not be tested if the new ones achieve the control objectives and ``have been in effect for a sufficient period to permit the auditor to assess their design and operating effectiveness'' (.53); the auditor may use the work of competent, objective others, less so as risk rises (.16 to .19); and AS 2601 governs service organizations (.B17). SEC guidance asks management for evidence ``from a reasonable period of time during the year, including the fiscal year-end'' \citep{sec2007guidance}. Neither sets a minimum period.

ITGCs over AI systems keep their four domains with new objects: change management over model versions, prompts, retrieval corpora, thresholds, tool permissions and vendor-initiated changes; logical access over agents' identities and write rights, separating configuration of the AI from approval of its outputs; operations that log prompts, outputs and model versions and monitor drift; and version control with pinned models, pre-promotion evaluation and rollback (Proposition~\ref{prop:ext-commonmode}(d)). COSO's 2026 guidance sets these expectations, treating models as configuration items \citep{coso2026genai}, under Principles 9 (changes that could significantly affect internal control) and 11 (general controls over technology) of the 2013 Framework \citep{coso2013framework}.

\begin{remark}[Reliance on AI-performed controls]\label{rem:acct-evidence}
AS 2201 allows benchmarking of ``entirely automated application controls'' unchanged since their baseline, given effective ITGCs over program changes, access and operations (.B28, .B29), and treats as suited to it controls that match a defined program in a stable application (.B31). In our judgment a control run by a model that its vendor updates, or whose output varies across runs, fails these factors: it is tested each year by re-performing or sampling outputs, like a manual control, and its outputs are information produced by the company (AS 1105.10) \citep{pcaob2026as1105}, while a deterministic rule applied to a model's proposal can still be benchmarked. A provider's SOC 2 report addresses the Trust Services Criteria; SOC 1 reports address controls relevant to user entities' ICFR \citep{aicpa2026soc1}. Banking model-risk guidance does not fill the gap: the agencies' revised guidance of 17 April 2026, which superseded SR 11-7, places generative and agentic AI models outside its scope and sets no enforceable standards \citep{fed2026sr262,occ2026bulletin13}. We found no PCAOB standard or staff guidance on AI-performed controls as of October 2026: the 2024 GenAI Spotlight reports outreach \citep{pcaob2024genai}, the 2025 inspection priorities say GenAI could amplify existing IT risks \citep{pcaob2024priorities}, and the 2024 technology-assisted analysis amendments, effective for fiscal years beginning on or after 15 December 2025, govern auditors' own procedures \citep{pcaob2024release007}. COSO expects controls relying on AI output to meet ICFR-level evidence standards, including a documented prompt, configuration and model version \citep{coso2026genai}.
\end{remark}

\begin{assumption}[Cutover of a financial-reporting process]\label{ass:acct-controls}
Process $p$ is in the scope of ICFR, fiscal years have unit length, and a cutover at position $z\in[0,1)$ leaves $H:=1-z$ to the year-end assessment.
\begin{enumerate}
\item[(i)] Control redesign costs $C^{\mathrm{ctl}}$ once (design, ITGCs over the AI system, testing, added audit work) and changes annual control costs by $\Delta C^{\mathrm{ctl}}$.
\item[(ii)] New controls count only after operating for $T_{\mathrm{op}}$ (a judgment parameter), and the close lasts $T_{\mathrm{cl}}$ after year-end; the firm cuts over only at $z\in[T_{\mathrm{cl}},1-T_{\mathrm{op}}]$, with $T_{\mathrm{op}}+T_{\mathrm{cl}}<1$.
\item[(iii)] With probability $p_{\mathrm{def}}$ the cutover creates a deficiency, a material weakness with probability $p_{\mathrm{mw}}$, detected after an exponential time with mean $T_{\mathrm{dc}}$; remediation takes $T_{\mathrm{rem}}$, after which the control must operate for $T_{\mathrm{op}}$. A material weakness unremediated at year-end costs $L_{\mathrm{mw}}$.
\item[(iv)] Delay costs $\hat c>0$ a year (forgone savings, or the parallel premium).
\end{enumerate}
\end{assumption}

\begin{proposition}[When to cut over a financial process]\label{prop:acct-cutover}
Under Assumption~\ref{ass:acct-controls}:
\begin{enumerate}
\item[(a)] (Calendar.) A share $1-T_{\mathrm{op}}-T_{\mathrm{cl}}$ of the year is open. If the ready date is uniform over the year and the firm cuts over at the first open date, the expected delay is $(T_{\mathrm{op}}+T_{\mathrm{cl}})^2/2$ years and its expected cost $\hat c(T_{\mathrm{op}}+T_{\mathrm{cl}})^2/2$.
\item[(b)] (Material-weakness risk.) The probability of an unremediated material weakness at the next year-end is
\begin{equation}\label{eq:acct-pmw}
P_{\mathrm{mw}}(H)=p_{\mathrm{def}}\,p_{\mathrm{mw}}\min\bigl\{1,\exp[-(H-T_{\mathrm{rem}}-T_{\mathrm{op}})/T_{\mathrm{dc}}]\bigr\},
\end{equation}
nonincreasing in $H$, and smallest among open dates right after the close, at $z=T_{\mathrm{cl}}$.
\item[(c)] (Cut over or wait.) A process ready at an open date $z$ should wait for the next close if and only if (without discounting)
\begin{equation}\label{eq:acct-defer}
L_{\mathrm{mw}}\bigl[P_{\mathrm{mw}}(1-z)-P_{\mathrm{mw}}(1-T_{\mathrm{cl}})\bigr]>\hat c\,(1+T_{\mathrm{cl}}-z).
\end{equation}
The difference between the sides is strictly increasing in $z$, so there is a threshold $z^*$: wait if and only if $z>z^*$.
\end{enumerate}
\end{proposition}

\begin{corollary}[Transition value for a financial process]\label{cor:acct-npv}
Under Assumptions~\ref{ass:tasks-transition} and \ref{ass:acct-controls}, with the controls redesigned at date 0, the cost change $\Delta C^{\mathrm{ctl}}$ on the success path after cutover, and a material weakness possible on every path that cuts over (valued at cutover),
\begin{equation}\label{eq:acct-npv}
\mathrm{NPV}^{\mathrm{ICFR}}_p=\mathrm{NPV}_p-C^{\mathrm{ctl}}-\varpi_p\Delta C^{\mathrm{ctl}}\frac{e^{-rT_{\mathrm{par}}}-e^{-rT_h}}r-[\varpi_p+(1-\varpi_p)p_u]\,e^{-rT_{\mathrm{par}}}L_{\mathrm{mw}}P_{\mathrm{mw}}(H),
\end{equation}
with $\mathrm{NPV}_p$ of Proposition~\ref{prop:tasks-transition}(a) evaluated at a parallel run that ends at an open date; the transition breakpoint rises by the last three terms.
\end{corollary}

\noindent\emph{Intuition and example.} The assessment is a year-end snapshot of controls that must have run long enough to be tested, so financial processes move just after the close and other processes fill the rest of the year. With $T_{\mathrm{op}}=0.25$ and $T_{\mathrm{cl}}=1/6$ (illustrative), \acctCutOpen{} of the year is open and the expected wait is \acctCutWait{} years, \acctCutCost{} at a delay cost of \acctCutDelay{} a year (illustrative; the invoice chain does not pay at the engine's defaults). With $p_{\mathrm{def}}=p_{\mathrm{mw}}=0.3$, a quarter to detect on average, a quarter to remediate and $L_{\mathrm{mw}}=\$5$ million, the risk of an unremediated material weakness is \acctCutPafter{} right after the close and \acctCutPlate{} with half a year or less to go, and a process ready after \acctCutZstar{} of the year should wait. On $L_{\mathrm{mw}}$ and $p_{\mathrm{def}}$: audit fees are about 35\% higher when deficiencies are disclosed \citep{hogan2008evidence}; the cost of equity moves by about 50 to 150 basis points after auditor-confirmed changes in control quality \citep{ashbaughskaife2009effect}; IT-related weaknesses come with more misstatements \citep{klamm2009sox}; and enterprise-system implementers later report fewer weaknesses \citep{morris2011impact}. Without a 404(b) audit, $L_{\mathrm{mw}}$ excludes an adverse auditor opinion, but management must still disclose a weakness it identifies. The trade in (c) is between delay and the chance that a deficiency introduced at cutover is still open at year-end; it does not license a design that leaves a material weakness, which Definition~\ref{def:acct-admissible} excludes.

\subsection{The accountability floor with AI-run controls}\label{sec:acct-amin}

\begin{assumption}[Layered controls]\label{ass:acct-layers}
Each in-scope control objective $o$ (completeness, accuracy, cutoff or valuation of a significant account, Definition~\ref{def:acct-scope}) has $N_o$ opportunities a year for a misstatement that would be material, and a design $d$ places layers of types $x$ (human; AI on a given provider) over it. A shared cause defeats every layer with probability $p^{\mathrm{cm}}_0$ per opportunity: management override, defects in source data or in information produced by the entity that every layer sees, or a failure of the IT general controls on which every layer's inputs depend. Otherwise layers of one type fail together with probability $p^{\mathrm{cm}}_x$ (common mode: shared error for people; a regression after a model update, drift or an injected instruction for a model) and independently with probability $p_x$, and types fail independently of one another. For an AI type, $p^{\mathrm{cm}}_x(\kappa_{\rm itgc})$ is nonincreasing in the quality $\kappa_{\rm itgc}$ of the ITGCs over its application (change management with pinned model versions and testing before promotion, logical access, operations), which are the controls over exactly these events. A layer of type $x$ costs $c_x$ a year; the application that runs the AI layers carries its ITGC line, $C^{\mathrm{itgc}}+C^{\mathrm{val}}/T_{\mathrm{up}}$ a year with $C^{\mathrm{val}}$ the retest at each forced upgrade every $T_{\mathrm{up}}$ years (one application; Proposition~\ref{prop:acct-itgc}(b) gives the general case); and the cost of accountability $C^{\mathrm{acct}}$ (certifications, audit committee, integrated audit, D\&O insurance) is common to all designs. An undetected material misstatement costs $L_{\mathrm{mm}}$ in consequences (restatement and remediation, higher audit fees, cost of capital, litigation, recovery of incentive pay), not the amount misstated.
\end{assumption}

The independence of types is an assumption, not a property of ICFR. Reviewers of AI output are subject to automation bias \citep{parasuraman2010complacency}, and models from different providers make correlated errors: on a benchmark of multiple-choice questions, when two models are both wrong they give the same wrong answer about 60\% of the time, against a third by chance \citep{kim2025correlated}. That figure is the probability of the same wrong answer given that both are wrong, not the probability that a second provider fails given that the first does, which a two-layer design needs and which is not measured. For a check that compares two providers and flags disagreement, an error escapes only if both err and agree, so the measured share is the escaping fraction of joint errors; the rate of joint errors stays a scenario value. Treating layers on different providers as independent types therefore gives an upper bound on the value of diversity, and the shared cause $p^{\mathrm{cm}}_0$ is where management override and common data enter.

\begin{definition}[Admissible control design]\label{def:acct-admissible}
Let $P_o(d)$ be the probability that a material misstatement at one opportunity of objective $o$ is not prevented or detected on a timely basis under design $d$, and $\Pi_o(d):=1-(1-P_o(d))^{N_o}$ the probability that at least one is missed in the year. A design is \emph{admissible} if and only if no in-scope control objective is left with a reasonable possibility that a material misstatement is not prevented or detected on a timely basis: $\Pi_o(d)\le\bar\epsilon$ for every in-scope $o$, where $\bar\epsilon$ is the probability up to which that possibility is remote. Otherwise the design has a material weakness (AS 2201, paragraph A7), and management may not conclude that ICFR is effective (Item 308(a)(3)) \citep{pcaob2026as2201,cfr2026sec229308}.
\end{definition}

A reasonable possibility is an event that is reasonably possible or probable in the sense of FASB Statement No.~5, so the complement is a remote event, one whose chance ``is slight'' \citep{fasb1975fas5}; and the severity of a deficiency depends on whether there is a reasonable possibility that controls fail to prevent or detect a misstatement, not on whether one occurs (AS 2201.63 and .64). No standard puts a number on $\bar\epsilon$, which is a judgment of management and the auditor; we report results at two illustrative values. The constraint is pass or fail: one material weakness bars an effective conclusion, so expected misstatement losses are not traded against control costs inside it.

\begin{proposition}[The accountability floor as an ICFR constraint]\label{prop:acct-amin}
Under Assumption~\ref{ass:acct-layers} and Definition~\ref{def:acct-admissible}, let a design use $n_x\ge 0$ layers of type $x$ over an objective, at least one in all, and write $\bar P_o:=1-(1-\bar\epsilon)^{1/N_o}$.
\begin{enumerate}
\item[(a)] (Depth.) A misstatement goes undetected with probability $P=p^{\mathrm{cm}}_0+(1-p^{\mathrm{cm}}_0)\prod_x[p^{\mathrm{cm}}_x+(1-p^{\mathrm{cm}}_x)p_x^{n_x}]$ (factor 1 for unused types). For one type, $P$ falls strictly with depth to $p^{\mathrm{cm}}_0+(1-p^{\mathrm{cm}}_0)p^{\mathrm{cm}}_x$; in general $P\ge p^{\mathrm{cm}}_0+(1-p^{\mathrm{cm}}_0)\prod_{x:n_x\ge1}p^{\mathrm{cm}}_x$.
\item[(b)] (Diversity.) Replacing one of $n$ layers of one AI type, with $P_n:=p^{\mathrm{cm}}+(1-p^{\mathrm{cm}})p^n$, by a layer of an independent type with failure probability $p^{\mathrm{ind}}$ lowers $P$ if and only if $p^{\mathrm{ind}}<P_n/P_{n-1}$, whatever $p^{\mathrm{cm}}_0$. This holds whenever $p^{\mathrm{ind}}<p$; and as $P_n/P_{n-1}\to 1$ with depth, a much less reliable independent layer beats another same-model layer once common-mode failure dominates.
\item[(c)] (Admissibility and the floor.) A design is admissible for $o$ if and only if $P_o(d)\le\bar P_o$. The accountability floor is the cost of the cheapest admissible design,
\begin{equation}\label{eq:acct-amin}
a_{\min}=\frac1S\min_{d\ \text{admissible}}\Bigl[C^{\mathrm{acct}}+\mathbf 1\{\text{AI used}\}\Bigl(C^{\mathrm{itgc}}+\frac{C^{\mathrm{val}}}{T_{\mathrm{up}}}\Bigr)+\sum_xn_xc_x\Bigr],
\end{equation}
summed over the objectives' designs, with no expected-loss term: an inadmissible design is excluded however cheap, and if no design is admissible the firm cannot assert effective ICFR at these parameters. A design whose types' common-mode floor $p^{\mathrm{cm}}_0+(1-p^{\mathrm{cm}}_0)\prod_{x:n_x\ge1}p^{\mathrm{cm}}_x$ exceeds $\bar P_o$ is inadmissible at every depth; in particular an all-AI design on one provider is admissible for some depth only if $p^{\mathrm{cm}}_0+(1-p^{\mathrm{cm}}_0)p^{\mathrm{cm}}<\bar P_o$, about $N_o(p^{\mathrm{cm}}_0+p^{\mathrm{cm}})<\bar\epsilon$.
\item[(d)] (Minimal depth.) For one type, the smallest admissible depth is $n_{\min}=\lceil\ln\rho/\ln p\rceil$ with $\rho:=[\bar P_o-p^{\mathrm{cm}}_0-(1-p^{\mathrm{cm}}_0)p^{\mathrm{cm}}]/[(1-p^{\mathrm{cm}}_0)(1-p^{\mathrm{cm}})]$ when $\rho>0$, and no depth is admissible otherwise. It is nondecreasing in $N_o$ and nonincreasing in $\bar\epsilon$. Depth addresses independent failures; only diversity and ITGC quality address common-mode ones.
\item[(e)] (ITGC quality.) If better ITGCs lower $p^{\mathrm{cm}}_x(\kappa_{\rm itgc})$ for the AI types, the admissible set grows with $\kappa_{\rm itgc}$ and the floor, minimized over designs and $\kappa_{\rm itgc}$ with the cost of the ITGCs included, never rises; holding $p^{\mathrm{cm}}_x$ fixed while charging for the ITGCs overstates it.
\item[(f)] (Scale.) With $N_o$ fixed, admissibility does not depend on revenue and $a_{\min}$ falls as $1/S$, but not to zero, since certifications are signed by people. If $N_o$ grew in proportion to revenue, $\bar P_o$ would tend to zero and every design with a positive common-mode floor would fail at some scale.
\end{enumerate}
\end{proposition}

\noindent\emph{Intuition and example.} Layers on one model share its failures, so depth buys little against them and diversity and ITGCs a great deal; and since the test is whether a material misstatement remains reasonably possible, a design that fails it is not a cheaper option but a material weakness. Part (f) reads $N_o$ as opportunities for misstatements material at the firm's own materiality, which rises with its size. Take a firm with \$1 billion of revenue and one objective with 24 opportunities a year for a material misstatement, human layers failing 5\% of the time at \$500{,}000 a year (common mode 0.2\%), AI layers failing 2\% at \$50{,}000 (common mode 0.5\% with effective ITGCs), no shared cause, and the AI layers run by an application with the invoice application's ITGC line of \acctAminLine{} a year (Proposition~\ref{prop:acct-itgc}); the other values are illustrative. No all-AI design on one provider is admissible at any depth or any remote threshold below \acctAminFloorAI: its common-mode term alone leaves an \acctAminFloorAI{} chance a year of missing a material misstatement. Human layers alone leave at least \acctAminFloorHu. At a remote threshold of 5\%, three human layers are admissible (\acctAminHumThreeMiss, \acctAminHumThreeCost{} a year), and the cheapest admissible design is one human layer over one AI layer, \acctAminFiveCost{} a year with an annual miss probability of \acctAminFiveMiss. At 1\%, no human-only design is admissible, and the cheapest is one human layer over two AI layers, \acctAminOneCost{} (\acctAminOneMiss). Neither threshold is a central value. The designs are conditional on the illustrative layer values: one human over one AI layer stays admissible at 5\% only while a shared cause $p^{\mathrm{cm}}_0$ stays below \acctAminPccFive{} per opportunity and while the human layer misses no more than \acctAminPhmFive{} of the misstatements the AI layer missed (automation bias \citep{parasuraman2010complacency} makes the two types dependent), and one over two at 1\% only while $p^{\mathrm{cm}}_0$ stays below \acctAminPccOne. If the ITGCs did not pin model versions or test before promotion, so that the AI common-mode probability were 2\% (illustrative), no design with one human layer would be admissible at 1\%, and the cheapest would be two human layers and one AI layer at \acctAminWeakOneCost: effective ITGCs lower the floor by \acctAminItgcValue{} a year, more than their line. The residual expected consequence of the chosen design, \acctAminFiveResid{} a year at 5\% and \acctAminOneResid{} at 1\% if each miss costs \$10 million, is not part of the floor; a firm that also weighs it adds layers above the floor (minimizing control cost plus expected consequences would choose one human and two AI layers at \acctAminLossCost{} in all), but no such weighing can admit a design that fails the test. Part I's statement that the floor may rise is the case in which AI layers on one provider, or weak ITGCs, make the cheaper designs inadmissible.

In Corollary~\ref{cor:ceiling}, then, $a_{\min}$ is the cost of the cheapest design that keeps every in-scope objective free of a material weakness: a fixed part ($C^{\mathrm{acct}}$, ITGC lines, retests and the layers the test requires) that falls with scale but is never zero, and a feasibility condition, set by common-mode failure and by the diversity and ITGC quality of the design rather than by headcount, under which some designs are not available at any price. The residual expected loss of the chosen design is an operating cost, not part of the floor. Proposition~\ref{prop:tasks-aggregate}(a), evaluated at a zero error rate, omits the feasibility condition; Proposition~\ref{prop:ext-commonmode} covers operational losses, not misstatements. Engine inputs are in \texttt{notes/accounting\_notes.md}.

\subsection{Scope, IT general controls and timeliness}\label{sec:acct-scope}

\paragraph{What the guidance says.} For the evaluation of ICFR, SEC guidance asks management to evaluate ``those IT general controls that are necessary for the proper and consistent operation of other controls designed to adequately address financial reporting risks'', in areas such as program development, program changes, computer operations and access to programs and data, and not those that primarily pertain to the efficiency or effectiveness of operations \citep{sec2007guidance}. AS 2201 makes the identification of IT risks and controls part of the top-down approach rather than a separate evaluation (.36, Note), lists a control's reliance on ITGCs among its risk factors and expects an automated control to be lower risk ``if relevant information technology general controls are effective'' (.47), and ties benchmarking to effective general controls and to the integrity of the files, tables, data and parameters an automated control uses (.B29, .B30). In our judgment a model's prompts, retrieval corpora, thresholds and tool permissions are such parameters. COSO's 2026 guidance does not repeat ITGC guidance (p.~6); it treats models and their configurations as configuration items subject to access control, change management and other IT operations controls and holds AI output that management relies on to ICFR evidence standards (p.~13), notes that vendor-driven changes may require notification obligations or independent verification (p.~11), and asks for periodic retesting without setting a frequency (p.~14) \citep{coso2026genai}. It does not address provider outages or manual fallback, service identities as such, or auditors' reliance, and no PCAOB standard or staff guidance addresses AI-performed controls (Remark~\ref{rem:acct-evidence}). The definition below is our scoping principle: scope is the processes, systems, integrations and policies that affect the timeliness, accuracy and completeness of financial reporting.

\begin{definition}[ICFR scope]\label{def:acct-scope}
Let $\mathcal Y^*$ be the in-scope systems (\texttt{in\_scope\_system} in the engine): those that initiate, authorize, record, process or report the transactions and data of significant accounts and disclosures, or on which a relevant control depends. A unit $u$ (activity, process step or application) is in scope, $u\in\mathcal U^*$, if at least one channel holds:
\begin{enumerate}
\item[(T)] (Transactions.) $u$ initiates, authorizes, records, processes or reports such transactions, or performs a relevant control;
\item[(S)] (Systems.) $u$ reads from or writes to a system in $\mathcal Y^*$;
\item[(I)] (Integrations.) $u$ creates or changes an integration that feeds a system in $\mathcal Y^*$ (\texttt{feeds\_in\_scope});
\item[(P)] (Policies.) $u$ executes, or is governed by, an ICFR-relevant policy (revenue recognition, credit and refund authority, delegation of authority, the close calendar).
\end{enumerate}
Reads count because an identity that can read an in-scope system is in its logical-access population, and its outputs are information produced by the company once a control uses them.
\end{definition}

\begin{assumption}[IT general controls and the close window]\label{ass:acct-itgc}
Let $\mathcal U^*_{\mathrm{AI}}$ be the AI applications in $\mathcal U^*$. Application $g$ carries the ITGC cost line
\begin{equation}\label{eq:acct-itgc}
C^{\mathrm{itgc}}_g=c^{\mathrm{acc}}\bigl(n^{\mathrm{id}}_g+n^{\mathrm{perm}}_g\bigr)+c^{\mathrm{chg}}n^{\mathrm{chg}}_g+C^{\mathrm{val}}_g/T_{\mathrm{up}}+c^{\mathrm{ops}}_g+c^{\mathrm{ifc}}n^{\mathrm{ifc}}_g
\end{equation}
a year: access over its $n^{\mathrm{id}}_g$ service identities and $n^{\mathrm{perm}}_g$ tool permissions; change management over $n^{\mathrm{chg}}_g$ model, prompt, tool and configuration changes a year; a retest $C^{\mathrm{val}}_g$ of the controls it performs or supports at each forced model upgrade, every $T_{\mathrm{up}}$ years, which prices the re-testing of their design and operating effectiveness, not the functional regression suite after the upgrade (the revalidation of Assumption~\ref{ass:tasks-ancillary}, which is in $W$) nor the upgrade's change ticket (one of the $n^{\mathrm{chg}}_g$ changes); operations $c^{\mathrm{ops}}_g$ (logging of prompts, outputs and model versions, monitoring, incident handling); and interface controls over its $n^{\mathrm{ifc}}_g$ integrations that feed in-scope systems. Their design is part of $C^{\mathrm{ctl}}$. Provider incidents that stall $g$ start in the close window, of length $T_{\mathrm{cw}}$ in each of $n_{\mathrm{cl}}$ closes a year, at rate $\nu_{\mathrm{out}}$, last an exponential time with mean $\bar D$, and delay the close by their duration at a cost $c_{\mathrm{del}}$ per unit of time; a delay beyond the slack $T_{\mathrm{sl}}$ misses the reporting deadline, at a cost $L_{\mathrm{late}}$. A manual fallback costs $C^{\mathrm{fb}}$ a year to maintain, takes over after $t_{\mathrm{sw}}\le T_{\mathrm{sl}}$, and costs $c_{\mathrm{fb}}$ per unit of time while it runs. The costs $c_{\mathrm{del}}$ and $c_{\mathrm{fb}}$ are in excess of the ordinary loss rate $x_p$ of Assumption~\ref{ass:ext-commonmode}.
\end{assumption}

\begin{proposition}[Scope, IT general controls and timeliness]\label{prop:acct-itgc}
Under Definition~\ref{def:acct-scope} and Assumption~\ref{ass:acct-itgc}:
\begin{enumerate}
\item[(a)] (Scope.) $\mathcal U^*$ is the union of the four channel sets, so it grows with $\mathcal Y^*$, with the integrations that feed it and with the ICFR-relevant policies; a unit outside (T) is in scope if and only if (S), (I) or (P) holds. An AI agent that performs a unit of $\mathcal U^*$, taking over its inputs, outputs, integrations and governing policies, meets at least that unit's channels: automation never removes a unit from scope, and the agent is in $\mathcal U^*_{\mathrm{AI}}$.
\item[(b)] (The floor.) Equation \eqref{eq:acct-amin} holds with $\mathbf 1\{\text{AI used}\}(C^{\mathrm{itgc}}+C^{\mathrm{val}}/T_{\mathrm{up}})$ replaced by $\sum_{g\in\mathcal U^*_{\mathrm{AI}}}C^{\mathrm{itgc}}_g+C^{\mathrm{tml}}$, with $C^{\mathrm{tml}}$ from (d). Lines are per application, not per layer: AI layers in one application share a line, layers in separate applications add one each, and an application in scope through (S), (I) or (P) adds its line though it performs no control. The lines do not change which designs are admissible, only what they cost. With one application per AI layer, each AI layer costs $c_x+C^{\mathrm{itgc}}_g$, and the cheapest admissible design uses no more AI layers (up to ties) than with one shared application, and no more as the line rises.
\item[(c)] (Transition costs.) In Corollary~\ref{cor:acct-npv}, $\Delta C^{\mathrm{ctl}}$ includes $\sum_{g\in\mathcal U^*_{\mathrm{AI},\mathrm{new}}}C^{\mathrm{itgc}}_g-\sum_{g\in\mathcal U^*_{\mathrm{AI},\mathrm{ret}}}C^{\mathrm{itgc}}_g$ over the applications added and retired, plus the change in $C^{\mathrm{tml}}$. The retests recur for as long as $g$ runs. If upgrade dates are stationary and independent of the cutover date, their expected present value over $[T_{\mathrm{par}},T_h]$ is that of the flow $C^{\mathrm{val}}_g/T_{\mathrm{up}}$; if the upgrade clock starts at cutover, it is
\begin{equation}\label{eq:acct-pvup}
\mathrm{PV}^{\mathrm{up}}_g=C^{\mathrm{val}}_g\,e^{-r(T_{\mathrm{par}}+T_{\mathrm{up}})}\,\frac{1-e^{-rn_{\mathrm{up}}T_{\mathrm{up}}}}{1-e^{-rT_{\mathrm{up}}}},\qquad n_{\mathrm{up}}=\bigl\lfloor(T_h-T_{\mathrm{par}})/T_{\mathrm{up}}\bigr\rfloor .
\end{equation}
\item[(d)] (Timeliness.) With $\nu_{\mathrm{cw}}:=n_{\mathrm{cl}}\nu_{\mathrm{out}}T_{\mathrm{cw}}$, close-window incidents cost $C^{\mathrm{tml}}_{\mathrm{wait}}=\nu_{\mathrm{cw}}[c_{\mathrm{del}}\bar D+L_{\mathrm{late}}e^{-T_{\mathrm{sl}}/\bar D}]$ a year without a fallback and $C^{\mathrm{tml}}_{\mathrm{fb}}=C^{\mathrm{fb}}+\nu_{\mathrm{cw}}\bar D[c_{\mathrm{del}}(1-e^{-t_{\mathrm{sw}}/\bar D})+c_{\mathrm{fb}}e^{-t_{\mathrm{sw}}/\bar D}]$ with one. The fallback pays if and only if $\Delta^{\mathrm{fb}}:=\bar De^{-t_{\mathrm{sw}}/\bar D}(c_{\mathrm{del}}-c_{\mathrm{fb}})+L_{\mathrm{late}}e^{-T_{\mathrm{sl}}/\bar D}>0$ and $\nu_{\mathrm{out}}>\nu^*:=C^{\mathrm{fb}}/(n_{\mathrm{cl}}T_{\mathrm{cw}}\Delta^{\mathrm{fb}})$. The timeliness cost is $C^{\mathrm{tml}}:=\min\{C^{\mathrm{tml}}_{\mathrm{wait}},C^{\mathrm{tml}}_{\mathrm{fb}}\}$.
\end{enumerate}
\end{proposition}

\noindent\emph{Intuition and example.} Scope follows the data, not the organization chart: support staff who issue refunds and credits written to billing are in scope through (S) and (P) though they record no accounting entry, and an agent that takes over the refunds through a new feed into billing adds (I). Automation adds applications to scope and removes none, and the ITGC line is its price. For an application performing the invoice chain's matching and exception steps, with three service identities and twelve tool permissions at \$2{,}000 each, 24 changes at \$2{,}500, a retest at each forced upgrade every \acctItgcTup{} years, \$55{,}000 of operations and \acctItgcNint{} integrations into in-scope systems at \$10{,}000 (the engine's unit costs; the counts and the operations cost are example settings, and $T_{\mathrm{up}}$ and the integrations are the engine's), the line is \acctItgcInv{} a year. We price the retest at the effort the 2025 SOX survey of KPMG reports, 16 hours per control tested, for about ten controls at a loaded \$150 an hour: \acctItgcVal{} a retest, or \acctItgcUp{} a year \citep{kpmg2025sox}. The same survey implies about \$57{,}500 a year of SOX program cost per in-scope system (\$2.3 million over 40 systems), of which about \$16{,}000 falls on ITGCs if cost follows the count of key controls (our division); this anchors a line of about \$80{,}000 a year per in-scope AI application, with a range of \$16{,}000 to \$350{,}000. An application at the engine's default counts that automates one in-scope step with one interface costs \acctItgcApp. The invoice application is above that anchor for reasons of its own: it changes twice as often as the default, feeds four in-scope systems and carries the top of the operations range, and its line sits \acctItgcPos{} of the way up the range. The survey's figure averages over in-scope systems that are mostly stable packaged applications and does not price change evidence, interface controls or monitoring separately, so it is a floor for a frequently changing AI application. The line lowers the invoice transition's risk-adjusted NPV by a further \acctItgcNPV, to \acctItgcNPVafter; the retests alone are worth \acctItgcPVup{} if the upgrade clock starts at cutover (\acctItgcPVflow{} with stationary dates). The refund agent's line is \acctItgcRef. In the example of Proposition~\ref{prop:acct-amin}, one application running both AI layers gives the floor of \acctAminOneCost{} at a remote threshold of 1\%; with one application per layer the admissible designs are the same, and the cheapest is still one human and two AI layers, at \acctAminSepOne. For timeliness, take four closes a year with ten-business-day windows, twelve provider incidents a year, two days of slack, \$50{,}000 a day of delay, \$2 million for a missed deadline and a \$100{,}000 fallback that takes over in 2.4 hours (illustrative): six-hour incidents cost \$25{,}000 a year and the fallback does not pay; day-long incidents cost \$616{,}000, mostly deadline risk, and the fallback cuts this to \$161{,}000, paying above 2.2 incidents a year. Each cost is counted once: $C^{\mathrm{tml}}$ holds only the close-window increment, the ordinary cost of outages stays in the common-mode term of $c_{\mathrm{err}}$ (Proposition~\ref{prop:ext-commonmode}(a)), the misstatement risk of a compressed close is in $P$, $C^{\mathrm{fb}}$ excludes the operations line $c^{\mathrm{ops}}_g$, and a forced upgrade appears in three lines that do not overlap: its functional regression suite in the revalidation of $W$, its change ticket in $c^{\mathrm{chg}}n^{\mathrm{chg}}_g$ and the re-testing of the controls it touches in $C^{\mathrm{val}}_g$. For the floor, each in-scope AI application with a line like the invoice application's adds about \acctItgcShare{} of revenue at \$1 billion. The direct cost of accountability at that size is about \acctAminDirect{} of revenue: the SOX program budget of a firm with less than \$5 billion of revenue in KPMG's survey, \acctAminKpmgSmall{} (read from a text extraction of the survey's chart, whose layout and labels were checked), and the average audit fee of accelerated filers, \acctAminAuditFee{} \citep{kpmg2025sox,auditanalytics2023fees}; with the audit fee at the S\&P 500 ratio, a lower bound for a filer of this size, it is \acctAminDirectLo. With the cheapest admissible design of Proposition~\ref{prop:acct-amin}, the bottom-up floor is about \acctItgcFloorFive{} at a remote threshold of 5\% and \acctItgcFloorOne{} at 1\%, and it reaches the lower end of the band in the calibration of Corollary~\ref{cor:ceiling}, 0.5\%, at about \acctItgcBandNFive{} and \acctItgcBandNOne{} such applications per \$1 billion of revenue (\acctItgcBandNLo{} with the lower audit fee), or \acctItgcBandNfbFive{} at 5\% if each needs the fallback; the favorable, central and conservative values are unchanged.
